
% LaTeX Beamer file — upgraded week15
% Morten Hjorth-Jensen, April 30, 2026

%-------------------- begin beamer-specific preamble ----------------------

\documentclass{beamer}

\usetheme{red_plain}
\usecolortheme{default}

\setbeamertemplate{navigation symbols}{}

\usepackage{pgf}
\usepackage{graphicx}
\usepackage{epsfig}
\usepackage{relsize}
\usepackage{fancybox}
\usepackage{fancyvrb}
\usepackage{minted}
\usepackage{amsmath,amssymb,bm}
\usepackage[T1]{fontenc}
\usepackage[utf8]{inputenc}
\usepackage{colortbl}
\usepackage[english]{babel}
\usepackage{tikz}
\usepackage{framed}
\usepackage{xcolor}
\usepackage{booktabs}

\beamertemplatetransparentcovereddynamic

\AtBeginSection[]
{
  \begin{frame}<beamer>[plain]
  \frametitle{}
  \tableofcontents[currentsection]
  \end{frame}
}

\newcommand{\shortinlinecomment}[3]{\note{\textbf{#1}: #2}}
\newcommand{\longinlinecomment}[3]{\shortinlinecomment{#1}{#2}{#3}}

\definecolor{linkcolor}{rgb}{0,0,0.4}
\hypersetup{
    colorlinks=true,
    linkcolor=linkcolor,
    urlcolor=linkcolor,
    pdfmenubar=true,
    pdftoolbar=true,
    bookmarksdepth=3
    }
\setlength{\parskip}{0pt}

\newenvironment{doconceexercise}{}{}
\newcounter{doconceexercisecounter}
\newenvironment{doconce:movie}{}{}
\newcounter{doconce:movie:counter}
\newcommand{\subex}[1]{\noindent\textbf{#1}}

% Convenience macros
\newcommand{\E}{\mathbb{E}}
\newcommand{\R}{\mathbb{R}}
\newcommand{\DKL}{D_{\mathrm{KL}}}
\newcommand{\Norm}{\mathcal{N}}
\newcommand{\ELBO}{\mathcal{L}}

\logo{{\tiny \copyright\ 1999-2026, Morten Hjorth-Jensen. Released under CC Attribution-NonCommercial 4.0 license}}

%-------------------- end beamer-specific preamble ----------------------

\raggedbottom
\makeindex

\begin{document}

\newcommand{\exercisesection}[1]{\subsection*{#1}}

\title{Advanced Machine Learning and Data Analysis for the Physical Sciences}
\author{Morten Hjorth-Jensen\inst{1}}
\institute{Department of Physics and Center for Computing in Science Education,
University of Oslo, Norway\inst{1}}
\date{April 30, 2026\\
{\tiny \copyright\ 1999-2026, Morten Hjorth-Jensen. Released under CC
Attribution-NonCommercial 4.0 license}}

\begin{frame}[plain,fragile]
\titlepage
\end{frame}

% --- TABLE OF CONTENTS -------------------------------------------------------
\begin{frame}[plain,fragile]
\frametitle{Plans for the week of April 27 -- May 1, 2026}
\begin{block}{Deep Generative Models: from VAEs to Diffusion}
\begin{enumerate}
  \item Reminder: Variational Autoencoders (VAEs) — key equations and limitations
  \item From VAEs to hierarchical VAEs and the diffusion limit
  \item Mathematical derivation of diffusion models (DDPM)
  \item The noise-prediction objective and connection to score matching
  \item Reverse diffusion and the denoising network
  \item DDPM PyTorch implementation on MNIST
  \item VAEs vs.\ Diffusion Models: pros, cons, and when to use each
\end{enumerate}
\end{block}
\end{frame}

% --- READINGS ----------------------------------------------------------------
\begin{frame}[plain,fragile]
\frametitle{Readings}
\begin{block}{Core references}
\begin{enumerate}
  \item Kingma \& Welling, \emph{Auto-Encoding Variational Bayes} (2014),
        \href{https://arxiv.org/abs/1312.6114}{arXiv:1312.6114}
  \item Sohl-Dickstein et al., \emph{Deep Unsupervised Learning using
        Nonequilibrium Thermodynamics} (2015),
        \href{https://arxiv.org/abs/1503.03585}{arXiv:1503.03585}
  \item Ho, Jain \& Abbeel, \emph{Denoising Diffusion Probabilistic Models}
        (2020), \href{https://arxiv.org/abs/2006.11239}{arXiv:2006.11239}
  \item Calvin Luo, \emph{Understanding Diffusion Models: A Unified Perspective},
        \href{https://arxiv.org/abs/2208.11970}{arXiv:2208.11970} (the main source
        for the derivations here)
  \item Kingma et al., \emph{Variational Diffusion Models},
        \href{https://arxiv.org/abs/2107.00630}{arXiv:2107.00630}
  \item Foster, \emph{Generative Deep Learning}, ch.\ 4 (VAEs) and ch.\ 8 (Diffusion)
\end{enumerate}
\end{block}
\end{frame}

% ===============================================================================
\section{Reminder: Variational Autoencoders}
% ===============================================================================

\begin{frame}[plain,fragile]
\frametitle{Generative models: the common thread}

All generative models we study share the same goal: learn a distribution
$p_{\bm{\theta}}(\bm{x})$ over data $\bm{x}$ such that we can both
\emph{evaluate} the density and \emph{sample} from it.

\begin{block}{Three major families (this course)}
\begin{tabular}{lll}
\toprule
\textbf{Model} & \textbf{Latent space} & \textbf{Key mechanism} \\
\midrule
Energy-based / RBM & fixed-dim $\bm{h}$ & Contrastive Divergence \\
VAE & fixed-dim $\bm{h}$ & ELBO maximization \\
\textbf{Diffusion} & \textbf{same dim as} $\bm{x}$ & \textbf{Iterative denoising} \\
\bottomrule
\end{tabular}
\end{block}

\vspace{2mm}
Today we derive diffusion models \emph{starting} from the VAE framework,
showing they can be understood as VAEs taken to a Markovian limit with
infinitely many latent hierarchies.
\end{frame}

% --- VAE ARCHITECTURE --------------------------------------------------------
\begin{frame}[plain,fragile]
\frametitle{VAE: Architecture Recap}

An autoencoder learns a mapping $\bm{x}\to\bm{h}\to\hat{\bm{x}}$ with
$\|\bm{x}-\hat{\bm{x}}\|^2$ as loss.  A \textbf{variational} autoencoder
replaces the deterministic encoder with a \emph{distribution}:

\begin{columns}
\column{0.5\textwidth}
\textbf{Encoder} (inference network):
\[
  q_{\bm{\phi}}(\bm{h}|\bm{x})
  = \Norm\!\bigl(\bm{h};\,\bm{\mu}_{\bm{\phi}}(\bm{x}),\,
    \bm{\sigma}^2_{\bm{\phi}}(\bm{x})\mathbf{I}\bigr)
\]
\textbf{Decoder} (generative network):
\[
  p_{\bm{\theta}}(\bm{x}|\bm{h})
\]
\textbf{Prior}:
\[
  p(\bm{h}) = \Norm(\bm{h};\bm{0},\mathbf{I})
\]
\column{0.5\textwidth}
\begin{block}{Goal}
Maximize the marginal log-likelihood
\[
  \log p_{\bm{\theta}}(\bm{x}) = \log\!\int p_{\bm{\theta}}(\bm{x},\bm{h})\,d\bm{h}
\]
which is intractable → variational lower bound.
\end{block}
\end{columns}
\end{frame}

% --- ELBO DERIVATION ---------------------------------------------------------
\begin{frame}[plain,fragile]
\frametitle{The ELBO derivation (from week 13)}

For any approximate posterior $q_{\bm{\phi}}(\bm{h}|\bm{x})$:
\begin{align*}
\log p(\bm{x})
&= \log \int \frac{p(\bm{x},\bm{h})}{q_{\bm{\phi}}(\bm{h}|\bm{x})}\,
             q_{\bm{\phi}}(\bm{h}|\bm{x})\,d\bm{h} \\[4pt]
&= \log \E_{q_{\bm{\phi}}(\bm{h}|\bm{x})}
   \!\left[\frac{p(\bm{x},\bm{h})}{q_{\bm{\phi}}(\bm{h}|\bm{x})}\right] \\[4pt]
&\geq \underbrace{\E_{q_{\bm{\phi}}(\bm{h}|\bm{x})}
   \!\left[\log\frac{p(\bm{x},\bm{h})}{q_{\bm{\phi}}(\bm{h}|\bm{x})}\right]}_{
   \ELBO(\bm{\phi},\bm{\theta})}
\quad\text{(Jensen's inequality)}
\end{align*}

Splitting the joint:
\[
\boxed{\ELBO(\bm{\phi},\bm{\theta})
= \underbrace{\E_{q_{\bm{\phi}}(\bm{h}|\bm{x})}\!\left[
  \log p_{\bm{\theta}}(\bm{x}|\bm{h})\right]}_{\text{Reconstruction}}
  -\underbrace{\DKL\!\bigl(q_{\bm{\phi}}(\bm{h}|\bm{x})\,\|\,p(\bm{h})\bigr)}_{\text{Regularisation}}}
\]
\end{frame}

% --- ELBO ANALYSIS -----------------------------------------------------------
\begin{frame}[plain,fragile]
\frametitle{VAE: ELBO interpretation}

\[
\log p(\bm{x}) = \underbrace{\ELBO(\bm{\phi},\bm{\theta})}_{\leq\,\log p(\bm{x})}
+ \underbrace{\DKL\!\bigl(q_{\bm{\phi}}(\bm{h}|\bm{x})\,\|\,
                          p(\bm{h}|\bm{x})\bigr)}_{\geq\,0}
\]

\begin{itemize}
  \item Since $\DKL\geq 0$, the ELBO is a \textbf{lower bound} on the log-evidence.
  \item Maximising the ELBO simultaneously:
        \begin{enumerate}
          \item maximises the reconstruction quality, and
          \item minimises the gap to the true posterior.
        \end{enumerate}
  \item The ELBO is tight ($= \log p(\bm{x})$) only when
        $q_{\bm{\phi}}(\bm{h}|\bm{x}) = p(\bm{h}|\bm{x})$.
\end{itemize}
\end{frame}

% --- REPARAMETERISATION ------------------------------------------------------
\begin{frame}[plain,fragile]
\frametitle{VAE: the reparameterisation trick}

To back-propagate through a \emph{sampling} node we write
\[
  \bm{h} = \bm{\mu}_{\bm{\phi}}(\bm{x})
           + \bm{\sigma}_{\bm{\phi}}(\bm{x})\odot\bm{\epsilon},
  \qquad \bm{\epsilon}\sim\Norm(\bm{0},\mathbf{I}).
\]

The gradient
$\nabla_{\bm{\phi}}\,\E_{q_{\bm{\phi}}(\bm{h}|\bm{x})}[f(\bm{h})]$
can now be estimated as
$\E_{\bm{\epsilon}\sim\Norm(\bm{0},\mathbf{I})}\!\left[\nabla_{\bm{\phi}}
  f\!\left(\bm{\mu}_{\bm{\phi}}(\bm{x})
    +\bm{\sigma}_{\bm{\phi}}(\bm{x})\odot\bm{\epsilon}\right)\right]$,
which is low-variance and unbiased.

\begin{block}{Closed-form KL for Gaussian encoder and prior}
When $q = \Norm(\bm{\mu}_{\bm{\phi}},\bm{\sigma}^2_{\bm{\phi}}\mathbf{I})$
and $p(\bm{h}) = \Norm(\bm{0},\mathbf{I})$:
\[
  \DKL(q\,\|\,p)
  = \frac{1}{2}\sum_{j=1}^{d_h}\!\left(
    \mu_j^2 + \sigma_j^2 - \log\sigma_j^2 - 1\right)
\]
This term regularises the latent code, preventing posterior collapse to a
Dirac delta.
\end{block}
\end{frame}

% --- VAE LIMITATIONS ---------------------------------------------------------
\begin{frame}[plain,fragile]
\frametitle{VAE: Strengths and Limitations}

\begin{columns}
\column{0.48\textwidth}
\begin{block}{Strengths}
\begin{itemize}
  \item Principled probabilistic framework
  \item Continuous, structured latent space
  \item Efficient single-pass inference
  \item Stable training (no adversary)
  \item Exact ELBO evaluation
  \item Interpretable disentangled representations
\end{itemize}
\end{block}
\column{0.48\textwidth}
\begin{block}{Limitations}
\begin{itemize}
  \item \textbf{Blurry samples} — Gaussian decoder averages over modes
  \item \textbf{Posterior collapse} — powerful decoder ignores $\bm{h}$
  \item Fixed-dim latent bottle-neck
  \item Gaussian prior may not fit complex data
  \item ELBO gap: $\log p(\bm{x})$ is never exactly recovered
  \item Limited expressiveness of mean-field approximation
\end{itemize}
\end{block}
\end{columns}
\vspace{2mm}
\textit{These limitations motivate going beyond a single VAE —
either by stacking them (hierarchical VAEs) or taking the diffusion limit.}
\end{frame}

% ===============================================================================
\section{From Hierarchical VAEs to Diffusion Models}
% ===============================================================================

\begin{frame}[plain,fragile]
\frametitle{Hierarchical (Markovian) VAEs}

We can generalise the single-latent VAE to a \textbf{chain of $T$ latent spaces}
$\bm{h}_1, \bm{h}_2,\ldots,\bm{h}_T$ with Markovian structure:

\begin{align*}
p(\bm{x}, \bm{h}_{1:T})
  &= p(\bm{h}_T)\,p_{\bm{\theta}}(\bm{x}|\bm{h}_1)
     \prod_{t=2}^{T}p_{\bm{\theta}}(\bm{h}_{t-1}|\bm{h}_t)
     \quad\text{(generative)}\\[4pt]
q_{\bm{\phi}}(\bm{h}_{1:T}|\bm{x})
  &= q_{\bm{\phi}}(\bm{h}_1|\bm{x})
     \prod_{t=2}^{T}q_{\bm{\phi}}(\bm{h}_t|\bm{h}_{t-1})
     \quad\text{(inference)}
\end{align*}

The ELBO for a Markovian hierarchical VAE becomes:
\[
\log p(\bm{x})
\geq
\E_{q_{\bm{\phi}}(\bm{h}_{1:T}|\bm{x})}
\!\left[\log\frac{p(\bm{h}_T)\,p_{\bm{\theta}}(\bm{x}|\bm{h}_1)
    \prod_{t=2}^{T}p_{\bm{\theta}}(\bm{h}_{t-1}|\bm{h}_t)}
  {q_{\bm{\phi}}(\bm{h}_1|\bm{x})\prod_{t=2}^{T}
    q_{\bm{\phi}}(\bm{h}_t|\bm{h}_{t-1})}\right]
\]
Each level adds a new reconstruction/regularisation trade-off.
\end{frame}

\begin{frame}[plain,fragile]
\frametitle{Three restrictions that give a VDM}

A \textbf{Variational Diffusion Model} (VDM) is a Markovian hierarchical VAE
subject to three additional constraints (Luo 2022):

\begin{block}{Constraint 1: Full-dimensional latents}
$\dim(\bm{h}_t) = \dim(\bm{x}_0)$ for all $t$.
The latent variable $\bm{h}_t \equiv \bm{x}_t$ lives in \emph{data space}, not a
compressed bottleneck.
\end{block}

\begin{block}{Constraint 2: Fixed (non-learned) Gaussian encoder}
\[
q(\bm{x}_t|\bm{x}_{t-1})
= \Norm\!\bigl(\bm{x}_t;\,\sqrt{\alpha_t}\,\bm{x}_{t-1},\,(1-\alpha_t)\mathbf{I}\bigr)
\]
The encoder is not trained; it is a fixed noise schedule.
\end{block}

\begin{block}{Constraint 3: Terminal distribution is standard Gaussian}
$\alpha_t$ is chosen so that as $t\to T$,
$q(\bm{x}_T|\bm{x}_0)\approx\Norm(\bm{0},\mathbf{I})$.
\end{block}
\end{frame}

% ===============================================================================
\section{Mathematics of Diffusion Models}
% ===============================================================================

\begin{frame}[plain,fragile]
\frametitle{The forward (noising) process $q$}

Define cumulative noise levels $\bar\alpha_t = \prod_{s=1}^{t}\alpha_s$.
By the Markov structure, there is a \emph{closed-form} for the marginal at
any time $t$ given only $\bm{x}_0$:

\begin{block}{Closed-form forward diffusion (key result)}
\[
\boxed{q(\bm{x}_t|\bm{x}_0)
= \Norm\!\bigl(\bm{x}_t;\,\sqrt{\bar\alpha_t}\,\bm{x}_0,\,
  (1-\bar\alpha_t)\mathbf{I}\bigr)}
\]
\end{block}

\textbf{Proof (induction):}
Starting from $q(\bm{x}_t|\bm{x}_{t-1})=\Norm(\sqrt{\alpha_t}\,\bm{x}_{t-1},(1-\alpha_t)\mathbf{I})$,
the Gaussian convolution formula gives:
\begin{align*}
  \bm{x}_t
  &= \sqrt{\alpha_t}\,\bm{x}_{t-1} + \sqrt{1-\alpha_t}\,\bm{\epsilon}_{t-1}\\
  &= \sqrt{\alpha_t\alpha_{t-1}}\,\bm{x}_{t-2}
     + \sqrt{\alpha_t(1-\alpha_{t-1})+1-\alpha_t}\,\bar{\bm{\epsilon}}\\
  &= \cdots = \sqrt{\bar\alpha_t}\,\bm{x}_0 + \sqrt{1-\bar\alpha_t}\,\bm{\epsilon}
\end{align*}
where $\bm{\epsilon}\sim\Norm(\bm{0},\mathbf{I})$.
\end{frame}

\begin{frame}[plain,fragile]
\frametitle{The forward process: noise schedule}

\begin{columns}
\column{0.55\textwidth}
Two common schedules for $\beta_t = 1-\alpha_t$:

\textbf{Linear} (Ho et al. 2020):
\[
  \beta_t = \beta_{\min} + \frac{t-1}{T-1}(\beta_{\max}-\beta_{\min})
\]
with $\beta_{\min}=10^{-4}$, $\beta_{\max}=0.02$, $T=1000$.

\textbf{Cosine} (Nichol \& Dhariwal 2021):
\[
  \bar\alpha_t = \frac{f(t)}{f(0)},\quad
  f(t) = \cos^2\!\!\left(\frac{t/T + s}{1+s}\cdot\frac{\pi}{2}\right)
\]
avoids near-zero SNR at the start.

\column{0.43\textwidth}
\begin{block}{Signal-to-noise ratio (SNR)}
\[
\mathrm{SNR}(t) = \frac{\bar\alpha_t}{1-\bar\alpha_t}
\]
$\mathrm{SNR}(0)\gg 1$ (clean), $\mathrm{SNR}(T)\approx 0$ (pure noise).
The VDM objective can be rewritten entirely in terms of SNR.
\end{block}
\end{columns}
\end{frame}

\begin{frame}[plain,fragile]
\frametitle{Posterior of forward process: $q(\bm{x}_{t-1}|\bm{x}_t,\bm{x}_0)$}

The reverse of the forward step conditioned on $\bm{x}_0$ is tractable:

\begin{block}{Forward posterior (exact)}
\[
q(\bm{x}_{t-1}|\bm{x}_t,\bm{x}_0)
= \Norm\!\bigl(\bm{x}_{t-1};\,\tilde{\bm{\mu}}_t(\bm{x}_t,\bm{x}_0),\,
  \tilde\beta_t\mathbf{I}\bigr)
\]
with
\begin{align*}
\tilde{\bm{\mu}}_t(\bm{x}_t,\bm{x}_0)
  &= \frac{\sqrt{\bar\alpha_{t-1}}\,\beta_t}{1-\bar\alpha_t}\,\bm{x}_0
   + \frac{\sqrt{\alpha_t}(1-\bar\alpha_{t-1})}{1-\bar\alpha_t}\,\bm{x}_t\\[4pt]
\tilde\beta_t
  &= \frac{1-\bar\alpha_{t-1}}{1-\bar\alpha_t}\,\beta_t
\end{align*}
\end{block}

This is the \emph{target} for the learned reverse process
$p_{\bm{\theta}}(\bm{x}_{t-1}|\bm{x}_t)$.
\end{frame}

\begin{frame}[plain,fragile]
\frametitle{Optimising the VDM: ELBO decomposition}

Substituting the VDM structure into the MHVAE ELBO and simplifying:

\begin{align*}
\log p(\bm{x}_0)
&\geq
\underbrace{\E_{q(\bm{x}_1|\bm{x}_0)}
  \!\left[\log p_{\bm{\theta}}(\bm{x}_0|\bm{x}_1)\right]}_{\mathcal{L}_0:\ \text{reconstruction}}\\
&\quad -\underbrace{\DKL\!\bigl(q(\bm{x}_T|\bm{x}_0)\,\|\,p(\bm{x}_T)\bigr)}_{\mathcal{L}_T:\ \text{prior matching, no params}}\\
&\quad -\underbrace{\sum_{t=2}^{T}\E_{q(\bm{x}_t|\bm{x}_0)}\!\left[
  \DKL\!\bigl(q(\bm{x}_{t-1}|\bm{x}_t,\bm{x}_0)\,\|\,
             p_{\bm{\theta}}(\bm{x}_{t-1}|\bm{x}_t)\bigr)
\right]}_{\mathcal{L}_{1:T-1}:\ \text{denoising matching terms}}
\end{align*}

$\mathcal{L}_T=0$ by design. $\mathcal{L}_0$ can be treated as a single
decoder step. The dominant cost is in $\mathcal{L}_{1:T-1}$.
\end{frame}

\begin{frame}[plain,fragile]
\frametitle{Reducing variance: conditioning on $\bm{x}_0$}

The \textbf{consistency term} in the first derivation involved
two random variables per step, giving high variance.

Conditioning on $\bm{x}_0$ at each step (using the Markov property):
\[
  q(\bm{x}_{t-1}|\bm{x}_t,\bm{x}_0)
  \;\;\text{replaces}\;\;
  q(\bm{x}_{t-1}|\bm{x}_t)
\]
reduces each KL term to an expectation over \emph{one} random variable
$\bm{x}_t \sim q(\bm{x}_t|\bm{x}_0)$, substantially lowering Monte Carlo variance.

\begin{block}{The improved ELBO}
Each denoising matching term is:
\[
  \mathcal{L}_{t-1}
  = \E_{q(\bm{x}_t|\bm{x}_0)}\!\left[
    \DKL\!\bigl(q(\bm{x}_{t-1}|\bm{x}_t,\bm{x}_0)\,\|\,
               p_{\bm{\theta}}(\bm{x}_{t-1}|\bm{x}_t)\bigr)
  \right]
\]
Since both distributions are Gaussian, the KL has a closed form.
\end{block}
\end{frame}

\begin{frame}[plain,fragile]
\frametitle{Parameterising the reverse process}

Choose $p_{\bm{\theta}}(\bm{x}_{t-1}|\bm{x}_t)
= \Norm(\bm{x}_{t-1};\bm{\mu}_{\bm{\theta}}(\bm{x}_t,t),\tilde\beta_t\mathbf{I})$.

The KL between two Gaussians with equal variance reduces to a squared
mean difference:
\[
  \mathcal{L}_{t-1}
  = \E_{q(\bm{x}_t|\bm{x}_0)}\!\left[
    \frac{1}{2\tilde\beta_t}\,
    \bigl\|\tilde{\bm{\mu}}_t(\bm{x}_t,\bm{x}_0)
          - \bm{\mu}_{\bm{\theta}}(\bm{x}_t,t)\bigr\|^2
  \right] + C
\]

We can further reparameterise $\tilde{\bm{\mu}}_t$ using
$\bm{x}_0 = \tfrac{1}{\sqrt{\bar\alpha_t}}\bigl(\bm{x}_t - \sqrt{1-\bar\alpha_t}\,\bm{\epsilon}\bigr)$:

\[
\tilde{\bm{\mu}}_t
= \frac{1}{\sqrt{\alpha_t}}\!\left(
    \bm{x}_t - \frac{\beta_t}{\sqrt{1-\bar\alpha_t}}\,\bm{\epsilon}
  \right)
\]

If the network predicts $\bm{\epsilon}_{\bm{\theta}}(\bm{x}_t,t)\approx\bm{\epsilon}$:
\[
\bm{\mu}_{\bm{\theta}}(\bm{x}_t,t)
= \frac{1}{\sqrt{\alpha_t}}\!\left(
    \bm{x}_t - \frac{\beta_t}{\sqrt{1-\bar\alpha_t}}\,\bm{\epsilon}_{\bm{\theta}}(\bm{x}_t,t)
  \right)
\]
\end{frame}

\begin{frame}[plain,fragile]
\frametitle{The noise-prediction objective (DDPM)}

Substituting into $\mathcal{L}_{t-1}$:
\[
  \mathcal{L}_{t-1}
  = \E_{\bm{x}_0,\bm{\epsilon}}\!\left[
    \frac{\beta_t^2}{2\tilde\beta_t\alpha_t(1-\bar\alpha_t)}\,
    \bigl\|\bm{\epsilon} - \bm{\epsilon}_{\bm{\theta}}(\bm{x}_t,t)\bigr\|^2
  \right]
\]

Ho et al.\ (2020) showed that \textbf{dropping the weight} and averaging
over $t$ gives a simpler, better-performing objective:

\begin{block}{DDPM training loss}
\[
\boxed{
\mathcal{L}_{\mathrm{simple}}
= \E_{t,\bm{x}_0,\bm{\epsilon}}\!\left[
  \bigl\|\bm{\epsilon}
         - \bm{\epsilon}_{\bm{\theta}}(\sqrt{\bar\alpha_t}\,\bm{x}_0
           + \sqrt{1-\bar\alpha_t}\,\bm{\epsilon},\;t)\bigr\|^2
\right]
}
\]
where $t\sim\mathrm{Uniform}[1,T]$, $\bm{\epsilon}\sim\Norm(\bm{0},\mathbf{I})$.
\end{block}

The network $\bm{\epsilon}_{\bm{\theta}}$ \textbf{predicts the noise} that was added.
\end{frame}

\begin{frame}[plain,fragile]
\frametitle{Connection to score matching}

The noise prediction objective has a deep connection to \textbf{score matching}.

The \emph{score} of a distribution is $\nabla_{\bm{x}}\log p(\bm{x})$.

For the forward-diffused data at time $t$:
\[
  \nabla_{\bm{x}_t}\log q(\bm{x}_t|\bm{x}_0)
  = -\frac{\bm{x}_t - \sqrt{\bar\alpha_t}\,\bm{x}_0}{1-\bar\alpha_t}
  = -\frac{\bm{\epsilon}}{\sqrt{1-\bar\alpha_t}}
\]

Therefore predicting $\bm{\epsilon}_{\bm{\theta}}(\bm{x}_t,t)\approx\bm{\epsilon}$
is equivalent to estimating the \emph{score}:
\[
  \bm{s}_{\bm{\theta}}(\bm{x}_t,t)
  = -\frac{\bm{\epsilon}_{\bm{\theta}}(\bm{x}_t,t)}{\sqrt{1-\bar\alpha_t}}
  \approx \nabla_{\bm{x}_t}\log q(\bm{x}_t|\bm{x}_0)
\]

This links DDPM to \textbf{score-based generative models}
(Song \& Ermon, 2019) and \textbf{stochastic differential equations}
(Song et al.\ 2021).
\end{frame}

\begin{frame}[plain,fragile]
\frametitle{The reverse sampling algorithm (DDPM)}

Given the trained denoiser $\bm{\epsilon}_{\bm{\theta}}$:

\begin{block}{Algorithm: DDPM reverse sampling}
\begin{enumerate}
  \item Sample $\bm{x}_T\sim\Norm(\bm{0},\mathbf{I})$.
  \item For $t = T, T-1, \ldots, 1$:
    \begin{enumerate}
      \item Sample $\bm{z}\sim\Norm(\bm{0},\mathbf{I})$ if $t>1$, else $\bm{z}=\bm{0}$.
      \item Compute $\bm{\mu}_{\bm{\theta}}(\bm{x}_t,t) =
            \frac{1}{\sqrt{\alpha_t}}\!\left(\bm{x}_t
            - \frac{\beta_t}{\sqrt{1-\bar\alpha_t}}
              \bm{\epsilon}_{\bm{\theta}}(\bm{x}_t,t)\right)$
      \item $\bm{x}_{t-1} = \bm{\mu}_{\bm{\theta}}(\bm{x}_t,t)
            + \sqrt{\tilde\beta_t}\,\bm{z}$
    \end{enumerate}
  \item Return $\bm{x}_0$.
\end{enumerate}
\end{block}

Each step is one neural network evaluation + cheap Gaussian update.
With $T=1000$ steps this is slow but produces high-quality samples.
\end{frame}

\begin{frame}[plain,fragile]
\frametitle{Three equivalent parameterisations}

The denoising network can predict three equivalent targets, each with
slightly different properties:

\begin{block}{}
\begin{center}
\begin{tabular}{lll}
\toprule
\textbf{Predict} & \textbf{Objective} & \textbf{Remark} \\
\midrule
$\bm{\epsilon}$ (noise) &
  $\|\bm{\epsilon} - \bm{\epsilon}_{\bm{\theta}}(\bm{x}_t,t)\|^2$ &
  DDPM; most common \\[4pt]
$\bm{x}_0$ (data) &
  $\|\bm{x}_0 - \hat{\bm{x}}_{\bm{\theta}}(\bm{x}_t,t)\|^2$ &
  Image prior; can over-smooth \\[4pt]
$\bm{v} = \sqrt{\bar\alpha_t}\,\bm{\epsilon}
  - \sqrt{1-\bar\alpha_t}\,\bm{x}_0$ (velocity) &
  $\|\bm{v} - \bm{v}_{\bm{\theta}}(\bm{x}_t,t)\|^2$ &
  Progressive Distillation \\
\bottomrule
\end{tabular}
\end{center}
\end{block}

The velocity parameterisation (Salimans \& Ho, 2022) is particularly
useful for few-step distillation and stable training at small $t$.
\end{frame}

% ===============================================================================
\section{DDPM PyTorch Implementation}
% ===============================================================================

\begin{frame}[plain,fragile]
\frametitle{PyTorch DDPM on MNIST — Overview}

This self-contained implementation covers:
\begin{enumerate}
  \item Imports and hyperparameters / noise schedule
  \item Data loading (MNIST, normalised to $[-1,1]$)
  \item Lightweight U-Net with sinusoidal time embedding
  \item Forward diffusion: \texttt{q\_sample}
  \item Training loss: MSE between predicted and true noise
  \item Training loop
  \item Reverse sampling: \texttt{p\_sample\_loop}
\end{enumerate}

\vspace{2mm}
\begin{block}{References}
Jackson-Kang's PyTorch diffusion tutorial
(\href{https://github.com/Jackson-Kang/Pytorch-Diffusion-Model-Tutorial}{GitHub})
and awjuliani's DDPM
(\href{https://github.com/awjuliani/pytorch-diffusion}{GitHub}).
\end{block}
\end{frame}

\begin{frame}[plain,fragile]
\frametitle{Imports and noise schedule}

\begin{minted}[fontsize=\fontsize{8pt}{8pt},linenos=false,baselinestretch=1.0,
               fontfamily=tt,xleftmargin=2mm]{python}
import torch, torch.nn as nn, torch.nn.functional as F
import torchvision
from torchvision import datasets, transforms
from torch.utils.data import DataLoader
import matplotlib.pyplot as plt, math

device = 'cuda' if torch.cuda.is_available() else 'cpu'

# --- Hyperparameters ---------------------------------------------------
T           = 300        # diffusion timesteps
beta_start  = 1e-4
beta_end    = 0.02
img_size    = 28
channels    = 1
batch_size  = 128
epochs      = 5
lr          = 2e-4

# --- Linear noise schedule ---------------------------------------------
betas           = torch.linspace(beta_start, beta_end, T, device=device)
alphas          = 1.0 - betas
alphas_cumprod  = torch.cumprod(alphas, dim=0)   # \bar{\alpha}_t
\end{minted}
\end{frame}

\begin{frame}[plain,fragile]
\frametitle{Sinusoidal time embedding}

A key ingredient is conditioning the denoiser on the timestep $t$.
We use sinusoidal embeddings (analogous to transformer positional encoding):

\begin{minted}[fontsize=\fontsize{8pt}{8pt},linenos=false,baselinestretch=1.0,
               fontfamily=tt,xleftmargin=2mm]{python}
class SinusoidalPosEmb(nn.Module):
    """Map scalar t in [0,1] to a d-dimensional embedding."""
    def __init__(self, dim):
        super().__init__()
        self.dim = dim

    def forward(self, t):          # t: (B,) in [0, 1]
        half = self.dim // 2
        freqs = torch.exp(
            -math.log(10000) *
            torch.arange(half, device=t.device) / (half - 1)
        )
        args  = t.unsqueeze(1) * freqs.unsqueeze(0)   # (B, half)
        emb   = torch.cat([args.sin(), args.cos()], dim=-1)
        return emb                                     # (B, dim)
\end{minted}
\end{frame}

\begin{frame}[plain,fragile]
\frametitle{Lightweight U-Net model}

\begin{minted}[fontsize=\fontsize{8pt}{8pt},linenos=false,baselinestretch=1.0,
               fontfamily=tt,xleftmargin=2mm]{python}
class SimpleUNet(nn.Module):
    def __init__(self, c, emb_dim=128):
        super().__init__()
        self.time_mlp = nn.Sequential(
            SinusoidalPosEmb(emb_dim),
            nn.Linear(emb_dim, emb_dim), nn.ReLU()
        )
        self.enc1 = nn.Conv2d(c,   64,  3, padding=1)
        self.enc2 = nn.Conv2d(64, 128,  3, padding=1)
        self.dec1 = nn.ConvTranspose2d(128, 64, 3, padding=1)
        self.dec2 = nn.ConvTranspose2d(64,  c,  3, padding=1)
        self.act  = nn.ReLU()
        self.proj = nn.Linear(emb_dim, 128)  # time -> channel

    def forward(self, x, t):   # x: (B,C,H,W), t: (B,) in [0,1]
        temb = self.proj(self.time_mlp(t)).view(-1, 128, 1, 1)
        h    = self.act(self.enc1(x))
        h    = self.act(self.enc2(h)) + temb
        h    = self.act(self.dec1(h))
        return self.dec2(h)
\end{minted}
\end{frame}

\begin{frame}[plain,fragile]
\frametitle{Forward diffusion and training loss}

\begin{minted}[fontsize=\fontsize{8pt}{8pt},linenos=false,baselinestretch=1.0,
               fontfamily=tt,xleftmargin=2mm]{python}
def q_sample(x0, t, noise=None):
    """
    Sample x_t = sqrt(alpha_bar_t) * x0 + sqrt(1-alpha_bar_t) * eps
    Implements the closed-form forward process q(x_t | x_0).
    """
    if noise is None:
        noise = torch.randn_like(x0)
    acp  = alphas_cumprod[t].view(-1, 1, 1, 1)
    return acp.sqrt() * x0 + (1 - acp).sqrt() * noise


def diffusion_loss(model, x0):
    """
    Sample t ~ Uniform[0, T-1], noise ~ N(0,I),
    corrupt x0 to x_t, predict noise, compute MSE.
    """
    B     = x0.size(0)
    t     = torch.randint(0, T, (B,), device=device)
    noise = torch.randn_like(x0)
    x_t   = q_sample(x0, t, noise)
    t_in  = t.float() / T              # normalise to [0,1]
    pred  = model(x_t, t_in)
    return F.mse_loss(pred, noise)
\end{minted}
\end{frame}

\begin{frame}[plain,fragile]
\frametitle{Training loop}

\begin{minted}[fontsize=\fontsize{8pt}{8pt},linenos=false,baselinestretch=1.0,
               fontfamily=tt,xleftmargin=2mm]{python}
transform = transforms.Compose([
    transforms.ToTensor(),
    transforms.Normalize((0.5,), (0.5,)),   # -> [-1, 1]
])
train_ds     = datasets.MNIST('.', train=True, download=True,
                              transform=transform)
train_loader = DataLoader(train_ds, batch_size=batch_size, shuffle=True)

model = SimpleUNet(channels).to(device)
opt   = torch.optim.Adam(model.parameters(), lr=lr)

for epoch in range(epochs):
    total = 0.0
    for x, _ in train_loader:
        x    = x.to(device)
        loss = diffusion_loss(model, x)
        opt.zero_grad(); loss.backward(); opt.step()
        total += loss.item()
    print(f"Epoch {epoch+1}/{epochs}  "
          f"loss={total/len(train_loader):.4f}")
\end{minted}
\end{frame}

\begin{frame}[plain,fragile]
\frametitle{Reverse diffusion: sampling}

\begin{minted}[fontsize=\fontsize{8pt}{8pt},linenos=false,baselinestretch=1.0,
               fontfamily=tt,xleftmargin=2mm]{python}
@torch.no_grad()
def p_sample_loop(model, shape):
    """DDPM reverse process: x_T -> x_0."""
    x = torch.randn(shape, device=device)   # start from pure noise
    for i in reversed(range(T)):
        t_in    = torch.full((shape[0],), i, device=device).float()/T
        eps_hat = model(x, t_in)            # predict noise

        beta_t = betas[i]
        acp_t  = alphas_cumprod[i]
        coef1  = 1.0 / alphas[i].sqrt()
        coef2  = beta_t / (1.0 - acp_t).sqrt()
        mu     = coef1 * (x - coef2 * eps_hat)

        if i > 0:
            sigma = beta_t.sqrt()
            mu   += sigma * torch.randn_like(x)
        x = mu
    return x

samples = p_sample_loop(model, (16, channels, img_size, img_size))
samples = samples.clamp(-1, 1).cpu()
grid    = torchvision.utils.make_grid(samples, nrow=4, normalize=True)
plt.imshow(grid.permute(1, 2, 0)); plt.axis('off'); plt.show()
\end{minted}
\end{frame}

% ===============================================================================
\section{Comparing VAEs and Diffusion Models}
% ===============================================================================

\begin{frame}[plain,fragile]
\frametitle{Mathematical parallel: VAE vs.\ Diffusion}

Both optimise a \textbf{variational lower bound (ELBO)} on $\log p(\bm{x})$.

\begin{center}
\begin{tabular}{lll}
\toprule
\textbf{Property} & \textbf{VAE} & \textbf{DDPM} \\
\midrule
Latent dim & $d_h \ll d_x$ (bottleneck) & $d_h = d_x$ (full dim) \\
\# latent levels & 1 & $T$ (e.g.\ 1000) \\
Encoder $q$ & \emph{learned} $q_{\bm{\phi}}(\bm{h}|\bm{x})$ & \emph{fixed} Gaussian \\
Inference & single forward pass & single forward pass \\
Decoder $p$ & \emph{learned} $p_{\bm{\theta}}(\bm{x}|\bm{h})$ & \emph{learned} denoiser \\
Sampling & single decoder call & $T$ denoiser calls \\
ELBO terms & 2 (recon + KL) & $T$ (sum of KL terms) \\
Training loss & ELBO (two terms) & noise-prediction MSE \\
Reparameterisation & $\bm{h}=\bm{\mu}+\bm{\sigma}\odot\bm{\epsilon}$ & $\bm{x}_t=\sqrt{\bar\alpha_t}\bm{x}_0+\sqrt{1-\bar\alpha_t}\bm{\epsilon}$ \\
\bottomrule
\end{tabular}
\end{center}
\end{frame}

\begin{frame}[plain,fragile]
\frametitle{VAE vs.\ Diffusion: Pros and Cons}

\begin{columns}
\column{0.5\textwidth}
\begin{block}{VAE}
\textcolor{green!60!black}{\textbf{Pros:}}
\begin{itemize}\small
  \item \textbf{Fast} inference and generation (1 pass each)
  \item \textbf{Compact} latent space; useful for downstream tasks (clustering, interpolation)
  \item \textbf{Stable} training with well-understood ELBO
  \item \textbf{Semantic} control via latent arithmetic
  \item Closed-form KL for Gaussian encoder
\end{itemize}
\textcolor{red!60!black}{\textbf{Cons:}}
\begin{itemize}\small
  \item \textbf{Blurry} samples (Gaussian decoder, $\ell_2$ averaging)
  \item \textbf{Posterior collapse} risk
  \item \textbf{Mode dropping} (ELBO gap)
  \item Mean-field approximation limits posterior quality
\end{itemize}
\end{block}

\column{0.5\textwidth}
\begin{block}{Diffusion (DDPM)}
\textcolor{green!60!black}{\textbf{Pros:}}
\begin{itemize}\small
  \item \textbf{Sharp, high-quality} samples (state of the art)
  \item \textbf{Stable} training (MSE on noise)
  \item \textbf{No mode collapse} — covers full data distribution
  \item Theoretically grounded via ELBO and score matching
  \item Flexible conditioning (text, class, image, audio)
\end{itemize}
\textcolor{red!60!black}{\textbf{Cons:}}
\begin{itemize}\small
  \item \textbf{Slow} sampling ($T$ = 100--1000 NFEs)
  \item \textbf{No compact} latent representation
  \item High memory/compute for large $T$
  \item Less natural for downstream representation learning
\end{itemize}
\end{block}
\end{columns}
\end{frame}

\begin{frame}[plain,fragile]
\frametitle{When to use which model}

\begin{block}{Use a VAE when\ldots}
\begin{itemize}
  \item You need a \emph{structured, compressed} latent space for downstream tasks
        (clustering, regression, anomaly detection, disentanglement).
  \item Real-time generation or inference is required.
  \item The data has low intrinsic dimensionality relative to observation space
        (e.g.\ molecular graphs, structured scientific data).
  \item Interpretable uncertainty quantification is important.
\end{itemize}
\end{block}

\begin{block}{Use a Diffusion model when\ldots}
\begin{itemize}
  \item \emph{Sample quality} is the primary goal (images, audio, video).
  \item Conditioning on complex, high-dimensional signals is needed
        (text-to-image, super-resolution, inpainting).
  \item Training stability and diversity are more important than speed.
  \item You have the compute budget for $T$-step sampling.
\end{itemize}
\end{block}
\end{frame}

\begin{frame}[plain,fragile]
\frametitle{The common ELBO: a unifying view}

Both models maximise a lower bound of the same form:
\[
  \log p(\bm{x}) \geq \ELBO(\bm{\phi},\bm{\theta})
  = \underbrace{\text{reconstruction likelihood}}_{\E[\log p_{\bm{\theta}}(\bm{x}|z)]}
  - \underbrace{\text{regularisation}}_{\DKL(q_{\bm{\phi}}\,\|\,p)}
\]

In a VAE, the KL is between encoder and prior over a \emph{single} latent $\bm{h}$.

In a DDPM, the sum $\sum_{t=1}^{T}\mathcal{L}_{t-1}$ plays the same role:
each term is a KL between the forward posterior $q(\bm{x}_{t-1}|\bm{x}_t,\bm{x}_0)$
and the learned reverse $p_{\bm{\theta}}(\bm{x}_{t-1}|\bm{x}_t)$, summed over
$T$ levels.

\begin{block}{Key insight}
Diffusion models are VAEs with:
(a) $T$ latent levels instead of 1,
(b) a fixed encoder instead of a learned one, and
(c) full-dimensional latents instead of a bottleneck.
\end{block}
\end{frame}

\begin{frame}[plain,fragile]
\frametitle{Techniques for faster diffusion sampling}

DDPM requires $T\approx 1000$ steps.  Several approaches accelerate this:

\begin{columns}
\column{0.5\textwidth}
\begin{block}{DDIM (Song et al.\ 2020)}
Replace the stochastic reverse step with a \emph{deterministic} ODE integration.
Same noise network, but 50--200$\times$ fewer function evaluations (NFEs)
for comparable quality.
\end{block}
\begin{block}{DPM-Solver (Lu et al.\ 2022)}
Analytical solution of the probability-flow ODE yields high-order methods.
As few as 10--20 NFEs for MNIST-quality images.
\end{block}

\column{0.5\textwidth}
\begin{block}{Latent Diffusion (Rombach et al.\ 2022)}
Run diffusion in a \emph{VAE latent space} instead of pixel space.
Combines VAE compression with diffusion quality:
lower dimensionality $\Rightarrow$ fewer NFEs needed.
\textbf{This is how Stable Diffusion works}.
\end{block}
\begin{block}{Consistency Models (Song et al.\ 2023)}
Distill the whole trajectory into a single-step generator;
matches diffusion quality with one or two NFEs.
\end{block}
\end{columns}
\end{frame}

\begin{frame}[plain,fragile]
\frametitle{Latent diffusion: VAE meets diffusion}

\textbf{Latent Diffusion Models} (LDMs) elegantly combine both worlds:

\begin{enumerate}
  \item Train a VAE: $\bm{x}\to\bm{z}\in\R^{d_z}$ with $d_z \ll d_x$.
  \item Run DDPM \emph{in the VAE latent space} $\bm{z}$.
  \item Decode with the VAE decoder: $\hat{\bm{z}}_0\to\hat{\bm{x}}$.
\end{enumerate}

\begin{block}{Why this works}
\begin{itemize}
  \item The VAE provides a perceptually compressed, semantic representation.
  \item Diffusion in latent space is much cheaper: $d_z = 4\times 32\times 32$
        vs.\ $d_x = 3\times 256\times 256$ for Stable Diffusion.
  \item All conditioning (text, depth, segmentation) is applied in latent space
        via cross-attention in the U-Net.
  \item The VAE's KL regularisation ensures the latent space is smooth enough
        for a diffusion model to learn.
\end{itemize}
\end{block}
\end{frame}

% ===============================================================================
\section{Summary and Outlook}
% ===============================================================================

\begin{frame}[plain,fragile]
\frametitle{Summary: VAEs and Diffusion — key equations}

\begin{center}\small
\begin{tabular}{lll}
\toprule
\textbf{Concept} & \textbf{VAE} & \textbf{DDPM} \\
\midrule
Encoder & $q_{\bm{\phi}}(\bm{h}|\bm{x})=\Norm(\bm{\mu}_{\bm{\phi}},\bm{\sigma}^2_{\bm{\phi}}\mathbf{I})$
        & $q(\bm{x}_t|\bm{x}_{t-1})=\Norm(\sqrt{\alpha_t}\bm{x}_{t-1},(1-\alpha_t)\mathbf{I})$ \\[2pt]
Reparam.\ & $\bm{h}=\bm{\mu}+\bm{\sigma}\odot\bm{\epsilon}$
          & $\bm{x}_t=\sqrt{\bar\alpha_t}\bm{x}_0+\sqrt{1-\bar\alpha_t}\bm{\epsilon}$ \\[2pt]
Prior     & $p(\bm{h})=\Norm(\bm{0},\mathbf{I})$
          & $p(\bm{x}_T)=\Norm(\bm{0},\mathbf{I})$ \\[2pt]
KL term   & $\DKL(q_{\bm{\phi}}\|\,p)$ (closed form)
          & $\sum_t\DKL(q(\bm{x}_{t-1}|\bm{x}_t,\bm{x}_0)\|\,p_{\bm{\theta}})$ \\[2pt]
Loss      & $-\ELBO = -\E[\log p]+\DKL$
          & $\E\|\bm{\epsilon}-\bm{\epsilon}_{\bm{\theta}}(\bm{x}_t,t)\|^2$ \\[2pt]
Sampling  & $\bm{h}\sim q$; one decoder call
          & $\bm{x}_T\sim\Norm$; $T$ reverse steps \\
\bottomrule
\end{tabular}
\end{center}
\end{frame}

\begin{frame}[plain,fragile]
\frametitle{Exercises}

\begin{block}{Mathematical exercises}
\begin{enumerate}
  \item Verify the closed-form marginal $q(\bm{x}_t|\bm{x}_0)$
        by induction from $q(\bm{x}_t|\bm{x}_{t-1})=\Norm(\sqrt{\alpha_t}\bm{x}_{t-1},(1-\alpha_t)\mathbf{I})$.
  \item Derive the forward posterior $q(\bm{x}_{t-1}|\bm{x}_t,\bm{x}_0)$
        using Bayes' rule and the Gaussian multiplication formula.
  \item Show that the denoising-matching ELBO term $\mathcal{L}_{t-1}$ reduces
        to the noise-prediction MSE when both distributions are Gaussian with
        equal variances.
  \item Compute the closed-form KL divergence for the VAE between
        $q_{\bm{\phi}}(\bm{h}|\bm{x})=\Norm(\bm{\mu},\mathrm{diag}(\bm{\sigma}^2))$
        and $p(\bm{h})=\Norm(\bm{0},\mathbf{I})$.
\end{enumerate}
\end{block}

\begin{block}{Implementation exercises}
\begin{enumerate}
  \item Replace the linear schedule with a cosine schedule and compare sample quality.
  \item Implement DDIM: replace the stochastic reverse step with the ODE update
        and measure speedup vs.\ quality trade-off.
  \item Train a VAE on MNIST and compare the blurriness of VAE samples with DDPM samples
        using FID or visual inspection.
\end{enumerate}
\end{block}
\end{frame}

\begin{frame}[plain,fragile]
\frametitle{Plans for next week}

\begin{block}{Next topics}
\begin{enumerate}
  \item Generative Adversarial Networks (GANs): architecture, training, mode collapse
  \item Connecting GANs, VAEs, and Diffusion: a unified perspective on generative models
  \item Flow-based models (normalizing flows): exact likelihood, invertible networks
  \item Score-based models and stochastic differential equations
  \item Discussion of Project 2
\end{enumerate}
\end{block}

\begin{block}{Key references for next week}
\begin{itemize}
  \item Goodfellow et al., \emph{Generative Adversarial Nets}, NIPS 2014
  \item Rezende \& Mohamed, \emph{Variational Inference with Normalizing Flows}, ICML 2015
  \item Song et al., \emph{Score-Based Generative Modeling through SDEs}, ICLR 2021
\end{itemize}
\end{block}
\end{frame}

\end{document}
